\documentclass{beamer}
\usepackage{amsfonts,amsmath,oldgerm,mathrsfs,booktabs}
\usepackage{subcaption}
\usepackage{natbib}
\usepackage{adjustbox} % 在导言区加
\usetheme{sintef}

%%%%%%%% usage of the package {bm} %%%%%%%%
\usepackage{bm}
% \mb for mathbf
\newcommand{\mb}[1]{{\mathbf{#1}}}
\newcommand{\mbh}[1]{{\hat{\mathbf{#1}}}}
\newcommand{\mbb}[1]{{\bar{\mathbf{#1}}}}
\newcommand{\mbu}[1]{{\ubar}{\mathbf{#1}}}

% for text
\newcommand{\txbf}[1]{{\textbf{#1}}}
\newcommand{\txit}[1]{{\textit{#1}}}
\newcommand{\txrd}[1]{{\textcolor{sintefred}{\textbf{#1}}}}

% underline + bold (safe)
\newcommand{\ubar}[1]{\underline{\bm{#1}}}
\newcommand{\bmu}[1]{\bm{#1}}          % 建议：向量就是粗体
\newcommand{\bmb}[1]{\bar{\bm{#1}}}
\newcommand{\bmh}[1]{\hat{\bm{#1}}}
\newcommand{\bmc}[1]{\check{\bm{#1}}}
\newcommand{\bmbh}[1]{\hat{\bar{\bm{#1}}}}
\newcommand{\bmbt}[1]{\tilde{\bar{\bm{#1}}}}
\newcommand{\bmt}[1]{\tilde{\bm{#1}}}
\newcommand{\bmo}[1]{\overline{\bm{#1}}}

% small math symbols (no mbox)
\newcommand{\nmbs}[1]{{\scriptsize $#1$}}
\newcommand{\nmbst}[1]{{\scriptsize $\tilde{#1}$}}
\newcommand{\nmbss}[1]{{\tiny $#1$}}

%%%%%%%% expectation and probability %%%%%%%%
\newcommand{\ep}[1]{\mathbb{E}_{\mathbb{P}}\!\left[#1\right]}
\newcommand{\ephat}[1]{\mathbb{E}_{\hat{\mathbb{P}}}\!\left[#1\right]}
\newcommand{\ephati}[1]{\mathbb{E}_{\hat{\mathbb{P}}_i}\!\left[#1\right]}
\newcommand{\epdag}[1]{\mathbb{E}_{\mathbb{P}^\dag}\!\left[#1\right]}
\newcommand{\epcheck}[1]{\mathbb{E}_{\check{\mathbb{P}}}\!\left[#1\right]}
\newcommand{\epstar}[1]{\mathbb{E}_{\mathbb{P}^\star}\!\left[#1\right]}

\newcommand{\pp}[1]{\mathbb{P}\!\left\{#1\right\}}
\newcommand{\ppstar}[1]{\mathbb{P}^\star\!\left\{#1\right\}}

\newcommand{\bbp}{\mathbb{P}}
\newcommand{\bbq}{\mathbb{Q}}
\newcommand{\bbr}{\mathbb{R}}

%%%%%%%% calligraphic sets %%%%%%%%
\newcommand{\cX}{\mathcal{X}}
\newcommand{\cU}{\mathcal{U}}
\newcommand{\cV}{\mathcal{V}}
\newcommand{\cZ}{\mathcal{Z}}
\newcommand{\testcolor}[1]{\colorbox{#1}{\textcolor{#1}{test}}~\texttt{#1}}

\usefonttheme[onlymath]{serif}

\titlebackground*{assets/background}

\newcommand{\hrefcol}[2]{\textcolor{cyan}{\href{#1}{#2}}}

\title{Problem Statement and Model Formulation}
% \subtitle{subtitle}
% \course{Master's Degree in Computer Science}
% \author{author1, author2}
\author{Peize Li, Nankai University
}
% \supervisor{}
% \IDnumber{1234567}
\date{\today}

\begin{document}

\maketitle

\section{Problem Setting}


%%%%%%%%%%%%%%%%%%%%%%%%%%%%%%%%%%%%%%%%%%%%%%%%%%%%%%%%%
% Page 1

\begin{frame}{Multi-location Newsvendor Problem}

Consider a company that manages inventory in a distribution system composed of $J$ nonidentical retailers. \\

\vspace{0.3cm}

\begin{itemize}
\item Random demand ${d}_j, \forall j\in[J]$, ${\bm d} = ({d}_j)_{j \in [J]}$.
\item Contextual information ${\bm x}_j \in \cX \subseteq \bbr^p$, such as promotions, weather and holidays. ${\mb X} = ({\bm x}_j)_{j \in [J]}$.
\item Random variable $( {\mb X}, {\bm d})$ with joint distribution $\bbp$
\item Conditional distribution $\bbp_{\bm d \mid \mb X}$ of $\bm d$ given $\mb X$.
\end{itemize}
\vspace{0.3cm}

During the selling season, demand $\bm d$ realizes and incurs the overage cost $h_j$ per unit unsold inventory and lost-sale penalty cost $b_j > 0$ per unit unmet demand for retailer $j$. The total cost is $\sum_{j \in [J]} C_j (z_j,  d_j)$, where $C_j (\bm z_j, \bm d_j)$ is:

{\large
\[
 C_j (z_j,  d_j) =  h_j (z_j - d_j)^+ + b_j (d_j -b_j)^+ 
\]
}
\end{frame}

% Page 2

\begin{frame}{Contextual Ctochastic Problem}

Suppose the planner observes covariate $\mb X$ that is predictive of demand before the decision. Then, we can solve the following contextual stochastic optimization problem to find the optimal order quantity $\bm z$ conditioned on observed covariate $\mb X$ \citep{Ban_Rudin_2019big}: 
{\large
\begin{align}
\min_{\bm z \geq \bm 0} \  \mathbb{E}_{\bbp_{\bm d \mid \mb X}} \left[\sum_{j \in [J]}   C_j (z_j,  d_j)  \mid \mb X \right] 
\label{problem: contextual stochastic optimization}
\end{align}
}
\vspace{0.3cm}
By linearity of expectation, we can decompose the expectation to:
{\large
\[
\min_{\bm z \geq \bm 0} \ \sum_{j \in [J]} \mathbb{E}_{\bbp^{(j)}_{\bm d \mid \mb X}} \left[  C_j (z_j,  d_j)  \mid \mb X \right]
\]
}
where $ {\bbp^{(j)}_{\bm d \mid \mb X}}$ denotes the marginal distribution of $d_j$ in $\bbp_{\bm d \mid \mb X}$. However, we do not have access to true distribution.

\end{frame}


% Page 3

\begin{frame}{Separate Data Sets (Multi-source Predit-then-Optimize)}

Instead of assuming an aggregated data set at disposal \citep{notz2022prescriptive}, we consider the setting that the decision maker has access to $J$ \alert{disjoint data sets collected by retailers individually}.

\begin{itemize}
    \item Each retailer $j$ is associated with a \txrd{separate data set} $S^{(j)}=\{({\bm x}_i^{(j)},{d}_i^{(j)})\}_{i = 1}^{N_j}$, where $({\bm x}_i^{(j)},{ d}_i^{(j)})$ is i.i.d. samples generated from a distinct distribution:
    {\large
    \[
    {\bm x}_i^{(j)} \sim \bbp^{(j)}_{\mb X}, { d}_i^{(j)} \mid {\bm x}_i^{(j)} \sim \bbp^{(j)}_{\bm d \mid \mb X}, \quad \forall j \in [J]
    \]
    }
    where $\bbp^{(j)}_{\mb X}$ denotes the distribution of $\bm x_j$, and $\bbp^{(j)}_{\bm d \mid \mb X}$ denotes the marginal distribution of ${\bm d \mid \mb X}$ for retailer $j$.
    \item $\mb D^{(j)} = ({ d}_{1}^{(j)}, \ldots, { d}_{N_j}^{(j)}) \in \bbr^{N_j}$, $\mb X^{(j)} = ({\bm x}_{1}^{(j)}, \ldots, {\bm x}_{N_j}^{(j)})^\top \in \bbr^{N_j \times p} $ and total sample size $N = \sum_{j \in [J]} N_j$.
    \item \txrd{Heterogenous Distibution}: $\bbp^{(j)}_{\mb X}$ and $\bbp^{(j)}_{\bm d \mid \mb X}$  can be heterogenous among retailers.
    \item \txrd{Local data set} $\cup_{j \in [J]} S^{(j)}$, \txrd{Global data set} $\cup_{j \in [J]} S^{(j)}$.
    \end{itemize}

\end{frame}


\begin{frame}{Separate Data Sets (Multi-source Predit-then-Optimize)}
\vspace{-0.3cm}
\begin{figure}
    \centering
    \includegraphics[width=1\textwidth]{aa.jpg}
\end{figure}


\end{frame}

% Page 4

\begin{frame}{Separate Data Sets (Multi-source Predit-then-Optimize)}

The decision is made in two steps: First, we estimate the conditional distribution $\bbp_{\bm d \mid \mb X}$ given covariate $\mb X$. Second, we adopt a DRO approach to find a decision immunized against distribution ambiguity.

\vspace{0.3cm}

To obtain an estimate of marginal distribution $\bbp^{(j)}_{\bm d \mid \mb X}$, we assume the demand $d$ and the covariates $\bm x$ of a single retailer is governed by:

\begin{itemize}
    \item \txrd{Global model}: shared across all retailers
    {\large
    \[
    d_{\text{global}}  = \bm x ^ \top {\bm \beta}_{\text{global}}  + \varepsilon_{\text{global}} 
    \]
    }
    \item \txrd{Local model}: specific to each retailer $j$
    {\large
    \[
    d_{\text{local}}^{(j)}  = \bm x ^ \top {\bm \beta}^{(j)}_{\text{local}}  + \varepsilon_{\text{local}}^{(j)}
    \]
    }
    \item \txrd{Assumption}: $x^1 =1$ to allow for an intercept term, ${\bm \beta}_{\text{global}}, {\bm \beta}^{(j)}_{\text{local}} \in \bbr^p$ are the nonzero linear regression coefficients with $\|{\bm \beta}_{\text{global}}  \|_2 \leq B$ and $\| {\bm \beta}^{(j)}_{\text{local}} \|_2 \leq B$, and $\varepsilon_{\text{global}}, \varepsilon_{\text{local}}^{(j)}$ are bounded random noise with zero mean, independent of covariate $\bm x$. Note that we do not impose that $\varepsilon_{\text{local}}^{(j)}$ is independent over $j \in [J]$.
    \end{itemize}

\end{frame}

% Page 4

\begin{frame}{Joint Estimator}

The global model and local models are \txrd{closely related}, in the sense that they follow similar demand pattern in a market.
{
\large
\[
{\bm \beta}^{(j)}_{\text{local}} =  {\bm \beta}_{\text{global}}  + \xi^{(j)}
\]
}
where $\xi^{(j)}$ is the \txrd{bias} between the global model and local model.  This relation indicates that we can leverage global data set $\cup_{j \in [J]} S^{(j)}$ to help predict the local demand $d_{\text{local}}^{(j)}$ with careful debias.

\vspace{0.3cm}

If the true parameters $ {\bm \beta}^{(j)}_{\text{local}}$ and ${\varepsilon}^{(j)}_{\text{local}}$ are known, then the conditional stochastic program \eqref{problem: contextual stochastic optimization} is equivalent to:

{
\large
\begin{equation}
\bm v^*({\mb X}) = \ 
\min_{ z_j \geq 0, j\in[J]}
 \sum_{j=1}^J 
\mathbb{E} \left[\,C_j \left( \bm z_j, \bm x_j^\top  {\bm \beta}^{(j)}_{\text{local}} + \varepsilon_{\text{local}}^{(j)} \right) \right]
\label{problem: CSP with regression}
\end{equation}
}
where the expectation is taken with respect to the distribution of $\varepsilon_{\text{local}}^{(j)}$. We cannot solve \eqref{problem: CSP with regression} directly because the $ {\bm \beta}^{(j)}_{\text{local}}$ and $\varepsilon_{\text{local}}^{(j)}$ are unknown.

\end{frame}


% Page 5

\begin{frame}{Joint Estimator}

Our goal is to estimate $ {\bm \beta}^{(j)}_{\text{local}}$ and $\varepsilon_{\text{local}}^{(j)}$ accurately from separate data sets $S^{(j)}$ for $j \in [J]$. Since we have a linear model, the \txrd{ordinary least square (OLS) estimator} is the most prevalent choice \citet{bastani2021predicting}. We make some assumptions to assure idenfication of regression coefficients:
\begin{itemize}
    \item \txrd{Standardization:}
    We standardize each column $\bm x_i^{(j)}$ of $\mb X^{(j)}$ such that
    $\|\bm x_i^{(j)}\|_2^2 = N_j$.
    We further define the \textit{sample covariance matrix} for each $j$ as
    \[
        \Sigma^{(j)} = \frac{1}{N_j}\, \mb X^{(j)\top} \mb X^{(j)} .
    \]
    The diagonal of $\Sigma^{(j)}$ satisfies
    $\operatorname{Diag}(\Sigma^{(j)}) = \bm 1 \in \mathbb{R}^p$. \\ This is a standard practice to normalize covariates when using regularized regression.

    \item \txrd{Positive definiteness:}
    The sample covariance matrix
    \(
        \Sigma^{(j)} = \frac{1}{N_j}\, \mb X^{(j)\top} \mb X^{(j)}
    \)
    is positive definite, with the smallest eigenvalue bounded below by
    $\psi > 0$. \\ The second assumption is standard in statistics for idenfication for local and global estimator.

    \item \txrd{Sparsity:}
    The bias vector $\xi^{(j)}$ is sparse; that is,
    $\|\xi^{(j)}\|_0 = s \ll p$ for each $j$. \\ The joint estimator of \citet{bastani2021predicting} is most effective when sparsity is present.
\end{itemize}

\end{frame}

% Page 6

\begin{frame}{Joint Estimator}

Our goal is to estimate $ {\bm \beta}^{(j)}_{\text{local}}$ and $\varepsilon_{\text{local}}^{(j)}$ accurately from separate data sets $S^{(j)}$ for $j \in [J]$. There are two natural way to do so:

\begin{itemize}
    \item \txrd{Local Regression}. To solve a regression problem for each local data set $S^{(j)}$. 
    {\large
    \[
        \hat{\bm \beta}_{\text{local}} = \arg \min_{\bm \beta} \left\{\frac{1}{N_j} \|\mb D^{(j)}  - \mb X^{(j) \top} {\bm \beta}  \|_2^2 \right\}
    \]}
    \txrd{Drawback:} local data sets often face data scarcity, which may lead to high-variance estimation. 
    \item \txrd{Global Regression}. To obtain a better generalization performance, an alternative way is to collaboratively train a single global estimator model on the global data set $\cup_{j \in [J]} S^{(j)}$. 
    {\large
    \[
    \hat{\bm \beta}_{\text{global}} = \arg \min_{\bm \beta} \left\{ \sum_{j \in [J]} \frac{1}{N_j} \|\mb D^{(j)}  - \mb X^{(j) \top} {\bm \beta}  \|_2^2 \right\}
    \]}
    \txrd{Drawback:} However, such a global model may perform well on average but generalize poorly for a particular retailer $j$, because there are distributional heterogeneity \citep{guo2025statistical, mo2024minimax}.
\end{itemize}
\end{frame}

% Page 7

\begin{frame}{Joint Estimator}

We adapt the joint estimator from \citet{bastani2021predicting} on the global and local data as follows:

\begin{itemize}
    \item \txrd{Step 1: Global Estimator}
    We run an OLS regression on the global data set
    $\bigcup_{j \in [J]} S^{(j)}$
    and obtain a \textit{global estimator}
    $\hat{\bm\beta}_{\text{global}}$ defined as
    \begin{equation}
        \hat{\bm\beta}_{\text{global}}
        =
        \arg\min_{\bm\beta}
        \left\{
        \sum_{j \in [J]}
        \frac{1}{N_j}
        \left\|
            \mb D^{(j)} - \mb X^{(j)\top}\bm\beta
        \right\|_2^2
        \right\}.
        \label{eq:global_estimator}
    \end{equation}

    \item \txrd{Step 2: Joint Estimator}
    For each data set $S^{(j)}$, we obtain a joint estimator
    $\hat{\bm\beta}_{\text{joint}}^{(j)}$
    via a LASSO regression.
    The regularization term is based on the deviation between
    the estimated coefficient
    $\hat{\bm\beta}_{\text{joint}}^{(j)}$
    and the global estimator
    $\hat{\bm\beta}_{\text{global}}$:
    \begin{equation}
        \hat{\bm\beta}_{\text{joint}}^{(j)}
        =
        \arg\min_{\bm\beta}
        \left\{
        \frac{1}{N_j}
        \left\|
            \mb D^{(j)} - \mb X^{(j)\top}\bm\beta
        \right\|_2^2
        +
        \lambda^{(j)}
        \left\|
            \bm\beta - \hat{\bm\beta}_{\text{global}}
        \right\|_1
        \right\}.
        \label{eq:joint_estimator}
    \end{equation}
    Here, $\lambda^{(j)} > 0$ is a \txrd{regularization parameter} that controls
    the influence of the global estimator
    $\hat{\bm\beta}_{\text{global}}$
    in the joint model.
    The \txrd{$\ell_1$-norm penalty exploits the sparsity structure} of the bias term
    $\xi^{(j)}$.
\end{itemize}

\end{frame}



% Page 8

\begin{frame}{Empirical Optimization Problem}

\begin{itemize}
    \item \txrd{Empirical Residual}: For a given ${\bm \beta}$, we define the following mapping $\hat{\varepsilon}({\bm \beta}, j, i)$:
    \[
    \hat{\varepsilon}({\bm \beta}, j, i)  = {d}_i^{(j)} - {\bm x}_i^{(j) \top} {\bm \beta}
    \]
    which show that empirical residual is dependent on the regression coefficient ${\bm \beta}$.
    \item \txrd{Predicted Demand Sample}: we also define the following mapping $\hat{d}$ to predict the demand at any retailer:
    \[
    \hat{d}({\bm \beta}, {\bm x}, j, i) = \bm x ^\top \bm \beta + \hat{\varepsilon}({\bm \beta}, j, i) = (\bm x - {\bm x}_i^{(j)})^ \top {\bm \beta} + {d}_i^{(j)}
    \]
    The predicted demand is obtained by adjusting the historical demand at ${\bm x}_i^{(j)}$ with the predicted demand change induced by covariate shift $(\bm x - {\bm x}_i^{(j)})$.
    \item \txrd{Predicted Empirical Distribution}: for given ${\bm \beta}$ and covariate ${\bm x}$, we can construct the predicted empirical distribution of local demand $d_j$, using empirical residuals on the data set $S^{(j)}$:
    \[
    \hat{\bbp}(\bm \beta,\bm x, j) = \frac{1}{N_j} \sum_{i=1}^{N_j} \delta_{\hat{d}({\bm \beta}, {\bm x}, j, i)}
    \]
    where $\delta_{\hat{d}({\bm \beta}, {\bm x}, j, i)}$ is a Dirac measture that places a unit mass at ${\hat{d}({\bm \beta}, {\bm x}, j, i)}$.
\end{itemize}
\end{frame}

% Page 9

\begin{frame}{Empirical Optimization Problem}

If all the joint estimators $\hat{\bm \beta}_{\text{joint}}^{(j)}$ are known, the empirical optimization problem corresponding to Problem \eqref{problem: CSP with regression} is defined as
{
\large
\begin{equation}
\begin{aligned}
\hat{\bm v}^{\text{ER}}({\mb X}) & = \ 
\min_{\bm z \geq \bm 0} \ \sum_{j \in [J]} \mathbb{E} \left[  C_j (z_j,  d_j)  \mid \mb X \right]
 \\
& = \ 
\min_{ z_j \geq 0, j\in[J]}
 \sum_{j=1}^J 
 \frac{1}{N_j} \sum_{i = 1}^{N_j} 
C_j \left( \bm z_j, \hat{d}(\hat{\bm \beta}_{\text{joint}}^{(j)}, \bm x_j, j, i)  \right) 
\label{problem: Empirical Optimization Problem}
\end{aligned}
\end{equation}
}
where the expectation is taken with respect to the predicted empirical distribution $\hat{\bbp}(\hat{\bm \beta}_{\text{joint}}^{(j)},\bm x_j, j)$ \citep{kannan2025data}. \\

\vspace{0.3cm}

\txrd{Limitation}: A limitation of the empirical optimization approach may exhibit disappointing out-of-sample performance when the sample size is limited \citep{kannan2024residuals}. To address such small-data issue, we develop a DRO model to leverage predicted empirical distributions from multiple data sets, ensuring generalization to the specific retailer $j$.

\end{frame}


% Page 10

\begin{frame}{Mixture DRO Approach}

We borrow ideas from the conditional group DRO framework and feaderate learning  \citep{marfoq2021federated, wang2025distributionally}, to define \txrd{the local distribution as mixtures from the predicted empirical distributions}. 

\vspace{0.3cm}
Based on input covariate $\mb X$, we define the mixture conditional distribution $\bbq^{(j)}_{\bm d \mid \mb X}$ as:
\[
\bbq^{(j)}_{\bm d \mid \mb X} = \sum_{l \in [J]} q_l^{(j)} \hat{\bbp}(\hat{\bm \beta}_{\text{joint}}^{(l)},\alert{\bm x_j}, l)
\]
where $\bm q^{(j)} \in \Delta^J$ and $\Delta^J$ denotes $(J-1)$-dimensional probability simplex.  We first \txrd{fix the input covariate $\mb X$}, and estimate the empirical distribution $\hat{\bbp}(\hat{\bm \beta}_{\text{joint}}^{(l)},\bm x_j, l)$ \txrd{using different data set $S^{l}$ for $l \in [J]$}.

\vspace{0.3cm}


\end{frame}

% Page 11

\begin{frame}{Mixture DRO Approach}

We borrow ideas from the conditional group DRO framework and feaderate learning  \citep{marfoq2021federated, wang2025distributionally}, to define \txrd{the local distribution as mixtures from the predicted empirical distributions}. 

\vspace{0.3cm}
We can construct an ambiguity set $d_j$ conditioning on the observed covariate $\mb X$:
{\large 
\[
\mathcal{F}_j(\mb X, \theta) = \left \{ 
	\bbq^{(j)}_{\bm d \mid \mb X} \in \mathcal{P}(\bbr_+) \ 
	\begin{array}{|ll} 
	\ d_j \sim \bbq^{(j)}_{\bm d \mid \mb X}	\\ 
	\
	\bbq^{(j)}_{\bm d \mid \mb X} = \sum_{l \in [J]} q_l^{(j)} \hat{\bbp}(\hat{\bm \beta}_{\text{joint}}^{(l)},\bm x_j, l)  \\
	\ \bm q^{(j)} \in \mathcal{H}(\hat{\bm q}^{(j)} , \theta)
	\end{array}
	\right  \}
\]
}
where the uncertainty set $\mathcal{H}(\hat{\bm q}, \theta) \subseteq \Delta^J$ is defined as:
{\large
\[
\mathcal{H}(\hat{\bm q}, \theta) = \left \{ \bm q \in \bbr^J \
\begin{array}{|ll} \
	\bm 1 ^\top \bm q = 1 \\
	\ \bm q \geq \bm 0 \\
	\ \|\bm q - \alert{\hat{\bm q}} \| \leq \alert{\theta}
\end{array}
\right \} 
\]
}
\begin{itemize}
    \item \txrd{Mixture weight: $\bm q$ reside in the $\theta$-neighborhood of $\hat{\bm q}$} with respect to the general norm. 
    \item \txrd{Prior knowledge:} \txrd{$\hat{\bm q}$ is a pre-specified weight} obtained from training. 
    \item \txrd{Robustness budget:} $\theta \geq 0$ controls the size of $\mathcal{H}(\hat{\bm q}, \theta)$.
\end{itemize}

\end{frame}

% Page 12

\begin{frame}{Mixture DRO Approach}

Then, the ambiguity set for the joint distribution of demand $\bm d$ given covariate $\mb X$ is:
{\large
\begin{equation}
\mathcal{F}(\mb X, \theta) = \left \{ 
	\bbq_{\bm d \mid \mb X} \in \mathcal{P}_0(\bbr^J_+) \ 
	\begin{array}{|ll} 
	\ (d_1, \ldots, d_J) \sim \bbq_{\bm d \mid \mb X} \\ 
	\ d_j \sim \bbq^{(j)}_{\bm d \mid \mb X} & \forall j \in [J] \\ \
	\alert{\bbq^{(j)}_{\bm d \mid \mb X} \in \mathcal{F}_j(\mb X, \theta)} & \forall j \in [J] 
	\end{array}
	\right  \}
\label{Conditioanl Ambiguity Set}
\tag{\sc{Mixture}}
\end{equation}
}

With the defined ambiguity set and observed covariate $\mb X$, we formulate the following contextual DRO model to minimize the worst-case expected total cost:
{\large
\begin{equation}
\hat{\bm v}^{\text{DRO}}({\mb X})  = \min_{\bm z \geq \bm 0} \sup_{\bbq_{\bm d \mid \mb X} \in \mathcal{F}(\mb X, \theta)}  \mathbb{E}_{\bbq_{\bm d \mid \mb X}} \left[\sum_{j \in [J]}   C_j (z_j,  d_j)  \mid \mb X \right] 
\label{Problem: DRO model}
\tag{\sc{DRO}}
\end{equation}
}
We also denote an optimal solution to Problem \eqref{Problem: DRO model} as $\hat{\bm z}^{\text{DRO}}({\mb X})$.

\end{frame}


% Page 14 
\begin{frame}{Problem Reformulation}

% There may exist closed-form solutions
We first consider the classic DRO formulation.
{\large
\begin{equation}
\hat{\bm v}^{\text{DRO}}({\mb X})  = \min_{\bm z \geq \bm 0} \sup_{\bbq_{\bm d \mid \mb X} \in \mathcal{F}(\mb X, \theta)}  \mathbb{E}_{\bbq_{\bm d \mid \mb X}} \left[\sum_{j \in [J]}   C_j (z_j,  d_j)  \mid \mb X \right] 
\tag{\sc{DRO}}
\end{equation}
}

Given $\mathcal{F}(\mb X, \theta)$ in \eqref{Conditioanl Ambiguity Set}, we have equivalently(from Proposition 1 in \citet{hu2025budget}):

\begin{equation}
\hat{\bm v}^{\mathrm{DRO}}(\mb X)
=
\left\{
\begin{aligned}
\min_{\bm\tau,\,\bm z} \quad
& \sum_{j \in [J]} \tau_j \\
\text{s.t.} \quad
& \sup_{\bbq^{(j)}_{\bm d \mid \mb X} \in \mathcal{F}_j(\mb X, \theta)}
\mathbb{E}_{\bbq^{(j)}_{\bm d \mid \mb X}}
\!\left[
    C_j(z_j, d_j) \mid \mb X
\right]
\le \tau_j,
\quad \forall j \in [J], \\
& \bm z \ge \bm 0,\quad \bm\tau \ge \bm 0 .
\end{aligned}
\right.
\end{equation}


\end{frame}

% Page 15
\begin{frame}{Problem Reformulation}

\begin{itemize}
    \item By linearity of expectation, we can write \eqref{Problem: DRO model} as:
    \[
    \hat{\bm v}^{\text{DRO}}({\mb X})  = \min_{\bm z \geq \bm 0} \sum_{j \in [J]} \sup_{\bbq_{\bm d \mid \mb X} \in \mathcal{F}(\mb X, \theta)}  \mathbb{E}_{\bbq_{\bm d \mid \mb X}} \left[  C_j (z_j,  d_j)  \mid \mb X \right] 
    \]
    \item We introduce epigraphical variable $\tau_j$ for all $j \in [J]$ to represent the worst-case expected cost for retailer $j$.
    \item  By definition of \eqref{Conditioanl Ambiguity Set} ambiguity set, for any $\bbq_{\bm d \mid \mb X} \in \mathcal{F}(\mb X, \theta)$, there exist $\bbq^{(j)}_{\bm d \mid \mb X} \in \mathcal{F}_j(\mb X, \theta) ,\forall j\in [J]$ to constitute the joint distribution $\bbq_{\bm d \mid \mb X} $.
    \[
    \sup_{\bbq_{\bm d \mid \mb X} \in \mathcal{F}(\mb X, \theta)}  \mathbb{E}_{\bbq_{\bm d \mid \mb X}} \left[ C_j (z_j,  d_j)  \mid \mb X \right] \leq \sup_{\bbq^{(j)}_{\bm d \mid \mb X} \in \mathcal{F}_j(\mb X, \theta)}  \mathbb{E}_{\bbq_{\bm d \mid \mb X}} \left[ C_j (z_j,  d_j)  \mid \mb X \right]
    \]
    \item Conversely, for any $\bbq^{(j)}_{\bm d \mid \mb X} \in \mathcal{F}_j(\mb X, \theta), \forall j\in [J]$, we can construct a joint distribution $\bbq_{\bm d \mid \mb X} \in \mathcal{F}(\mb X, \theta)$ so that:
    \[
    \sup_{\bbq_{\bm d \mid \mb X} \in \mathcal{F}(\mb X, \theta)}  \mathbb{E}_{\bbq_{\bm d \mid \mb X}} \left[ C_j (z_j,  d_j)  \mid \mb X \right] \leq \sup_{\bbq^{(j)}_{\bm d \mid \mb X} \in \mathcal{F}_j(\mb X, \theta)}  \mathbb{E}_{\bbq_{\bm d \mid \mb X}} \left[ C_j (z_j,  d_j)  \mid \mb X \right]
    \]
    This completes the proof.
\end{itemize}

% There may exist closed-form solutions

\end{frame}


% Page 17
\begin{frame}{Problem Reformulation}
We next aim to transform the Problem (6) into a convex optimization problem. % There may exist closed-form solutions
We consider the reformulation of the following constraint:
\begin{equation}
\sup_{\bbq^{(j)}_{\bm d \mid \mb X} \in \mathcal{F}_j(\mb X, \theta)}  \mathbb{E}_{\bbq^{(j)}_{\bm d \mid \mb X}} \left[ C_j (z_j,  d_j)  \mid \mb X \right] \leq \tau_j
\label{eq: formulated constraint}
\end{equation}

For notational simplicity, we drop the subscript $j$ and write $\hat{\bbp}^{(j)}_l = \hat{\bbp}(\hat{\bm \beta}_{\text{joint}}^{(l)},\bm x_j, l)$. By definition,\txrd{$\bbq_{\bm d \mid \mb X}$ is a mixture distribution of $\hat{\bbp}(\hat{\bm \beta}_{\text{joint}}^{(l)},\bm x, l)$}.
\[
\mathbb{E}_{\bbq_{\bm d \mid \mb X}} \left[ C (\bm z,  d)  \mid \mb X \right] = \sum_{l \in [J]}  {q}_l \mathbb{E}_{\hat{\bbp}_l}[C (\bm z,  d)]
\]
Then the left-hand term can be simplified as
\[
\left\{
\begin{aligned}
\sup_{\bm q} \quad
& \sum_{l \in [J]} q_l \,
\mathbb{E}_{\hat{\bbp}_l}\!\left[ C(\bm z, d) \right] \\
\text{s.t.} \quad
& \bm 1^\top \bm q = 1, \\
& \|\bm q - \hat{\bm q}\| \le \theta, \\
& \bm q \ge \bm 0 .
\end{aligned}
\right.
\]
\end{frame}

% Page 18
\begin{frame}{Problem Reformulation}
By the definition of dual norm, we substitute the norm constraint by the following:
\[
\bm y^\top (\bm q - \hat{\bm q}) \le \theta,
\quad \forall\, \bm y \text{ such that } \|\bm y\|_* \le 1 .
\]
Then we can write the Lagrangian of the problem by introducing dual variable $\lambda \in \bbr$ and $\bm \nu \in \bbr^J_+$:
	\[
	L(\bm q, \lambda, \bm \nu, \bm y)= \mb E ^\top \bm q + \lambda (1 -  \bm 1^\top \bm q) + \theta - \bm y ^ \top (\bm q - \hat{\bm q})  + \bm \nu^\top \bm q 
	\]
where \txrd{we use $ \mb E ^\top \bm q $ to denote $\sum_{l \in [J]}  {q}_l \mathbb{E}_{\hat{\bbp}^{}_l}[C (\bm z,  d)]$} for short. By rearranging the terms, we have:
\[
	L(\bm q, \lambda, \bm \nu, \bm y) = (\mb E - \lambda \bm 1 - \bm y + \bm \nu )^\top \bm q + \lambda + \theta - \bm y ^\top \hat{\bm q}
\]
The result dual problem is:
\[
\left\{
\begin{aligned}
\inf_{\lambda,\, \bm\nu \ge \bm 0} \quad
& \lambda + \theta
- (\mb E - \lambda \bm 1 + \bm\nu)^\top \hat{\bm q}
 \\
\text{s.t.} \quad
& \|\mb E - \lambda \bm 1 + \bm\nu\|_* \le 1 .
\end{aligned}
\right.
\]

\end{frame}


% Page 19
\begin{frame}{Problem Reformulation}
Given $\mathcal{F}(\mb X, \theta)$ in \eqref{Conditioanl Ambiguity Set},
the Problem \eqref{Problem: DRO model} is equivalent to the following convex optimization problem:
\begin{equation}
\hat{\bm v}^{\mathrm{DRO}}(\mb X)
=
\left\{
\begin{aligned}
\min_{\bm\tau,\, \bm z,\, \{\lambda_j, \bm\nu_j\}_{j=1}^J} \quad
& \sum_{j \in [J]} \tau_j \\
\text{s.t.} \quad
& \lambda_j
- (\mb E_j - \lambda_j \bm 1 + \bm\nu_j)^\top \hat{\bm q}
+ \theta
\le \tau_j,
\quad \forall j \in [J], \\
& \|\mb E_j - \lambda_j \bm 1 + \bm\nu_j\|_* \le 1,
\quad \forall j \in [J], \\
& \lambda_j \in \mathbb{R}, \quad \bm\nu_j \ge \bm 0,
\quad \forall j \in [J], \\
& \bm z \ge \bm 0, \quad \bm\tau \ge \bm 0 .
\end{aligned}
\right.
\label{Problem:SummationDRO}
\end{equation}

\begin{itemize}
    \item \txrd{Advantage of Empirical distribution}: Note that we write  $ \mb E_j =(\mathbb{E}_{\hat{\bbp}^{(j)}_l}[C_j (z_j,  d_j)])_{l\in [J]}$. The expectation over empirical distribution is computationally tractable:
    \[
    \mathbb{E}_{\hat{\bbp}^{(j)}_l}[C_j (z_j,  d_j)] = \frac{1}{N_j}[C_j (z_j,  d_j)] = \frac{1}{N_j} \sum_{i = 1}^{N_j} 
    C_j \left( \bm z_j, \hat{d}(\hat{\bm \beta}_{\text{joint}}^{(j)}, \bm x_j, j, i) \right)
    \] 
\end{itemize}

\end{frame}


\begin{frame}{Potential Outcomes}

\begin{itemize}
    \item \txrd{Conjecture: $z^*$ for Problem \eqref{Problem: DRO model} for any retailer is a weighted average of individual empirical newsvendor solution}. 
    There may exist \txrd{closed-form solution} that for some $j \in [J]$, $\bm q^* \in \mathcal{H}(\hat{\bm q}, \theta)$:
    \[
    z^* = \arg \min_z \sum_{l \in [J]} {q}_l^* \mathbb{E}_{\hat{\bbp}_l}[C (\bm z, d)] = \sum_{l \in [J]} {q}_l^* z^{(l)}
    \]
    where \txrd{$z^{(l)}$ is empirical newsvendor solution for $l-$th SAA}:
    \[
    z^{(l)} = \arg \min_z \mathbb{E}_{\hat{\bbp}_l}[C (\bm z, d)] 
     = \inf\left\{z:\hat{F}_{l}(z)\geq\frac{b}{b+h}\right\}
    \]
    $\hat{F}_{l}$ is the empirical cumulative distribution
    function (cdf) of the demand from the $l-th$ dataset.
\end{itemize}


\end{frame}


% Page 13
\begin{frame}{Mixture DRO Approach}

\begin{itemize}
    \item \txrd{Extension 1: Budget-driven DRO} \citep{hu2025budget}. Given a budget $\tau_j$ for each $j$ and maximize the robustness budget $\theta$.
    \[
    \theta^*(\bm\tau)
    =
    \left\{
    \begin{aligned}
    \sup_{\theta} \quad & \theta \\
    \text{s.t.} \quad
    & \sup_{\substack{
    \bbq_{\bm d \mid \mb X} \in \mathcal{F}(\mb X, \theta)
    }}
    \mathbb{E}_{\bbq_{\bm d \mid \mb X}}
    \!\left[
        C_j(\bm z_j, d_j) \mid \mb X
    \right]
    \leq \tau_j,
    \quad \forall j \in [J], \\
    & \theta \ge 0 .
    \end{aligned}
    \right.
    \]

    \item \txrd{Extension 2: Robust Satisficing} \citep{sim2024analytics}. Given a budget $\tau_j$ for each $j$, and minimize the fragility measure. We can further account for the uncertainty in regression coefficients $\bm \beta$. 
    \[
    \kappa^*(\bm\tau)
    =
    \left\{
    \begin{aligned}
    \min_{\kappa} \quad & \kappa \\
    \text{s.t.} \quad
    & \mathbb{E}_{\bbq_{\bm d \mid \mb X}}
    \!\left[
        C_j(\bm z_j, d_j) \mid \mb X
    \right] - \tau_j
    \leq \kappa (\|\bm q - \hat{\bm q}\| - \theta)^+  ,
    \quad \forall  \bbq_{\bm d \mid \mb X}, j \in [J], \\
    & \kappa \ge 0 .
    \end{aligned}
    \right.
    \]

    \item \txrd{Extension 3: Minimum volume confidence set} \citep{lou2024estimation}:  $\mathcal{H}(\hat{\bm q}, \theta) \Rightarrow  \mathcal{H}(\hat{\bm q}, \mb D, \theta)$.
    \[
    \mathcal{H}(\hat{\bm q}, \mb D, \theta) =
    \mathcal{H}(\hat{\bm q}, \theta) = \left \{ \bm q \in \Delta^L \
    \begin{array}{|ll} \
        \ \|\alert{\mb D^-} (\bm q - \alert{\hat{\bm q}}) \| \leq \alert{\theta}
    \end{array}
    \right \} 
    \]

\end{itemize}

\end{frame}

%Page20

\begin{frame}{Mixture DRO Approach}

\begin{itemize}
    \item \txrd{Extension 3: Minimum volume confidence set (MVCS)} \citep{lou2024estimation}:  $\mathcal{H}(\hat{\bm q}, \theta) \Rightarrow  \mathcal{H}(\hat{\bm q}, \mb D, \theta)$.
    \[
    \mathcal{H}(\hat{\bm q}, \mb D, \theta) =
    \mathcal{H}(\hat{\bm q}, \theta) = \left \{ \bm q \in \Delta^L \
    \begin{array}{|ll} \
        \ \|\alert{\mb D^-} (\bm q - \alert{\hat{\bm q}}) \| \leq \alert{\theta}
    \end{array}
    \right \} 
    \]
    Similar to \citet{lou2024estimation}, we can construct MVCS problem:
        \[
    \mb D^*(\bm\tau)
    =
    \left\{
    \begin{aligned}
    \inf \quad & {\det} (\mb D) \\
    \text{s.t.} \quad
    & \sup_{\substack{
    \bbq_{\bm d \mid \mb X} \in \mathcal{F}(\mb X, \theta)
    }}
    \mathbb{E}_{\bbq_{\bm d \mid \mb X}}
    \!\left[
        C_j(\bm z_j, d_j) \mid \mb X
    \right]
    \leq \tau_j,
    \quad \forall j \in [J], \\
    & \mathbf{D}\succeq \boldsymbol{r}
    \end{aligned}
    \right.
    \]
    

\end{itemize}

\end{frame}


%%%%%%%%%%%%%%%%%%%%%%%%%%%%%%%%%%%%%%%%%%%%%%%%%%%%%%%%%
%%%%%%%%%%%%%%%%%%%%%%%%%%%%%%%%%%%%%%%%%%%%%%%%%%%%%%%%%

\begin{frame}[allowframebreaks]{References}
\tiny
\bibliographystyle{apalike}
\bibliography{bibliography}
\end{frame}

\backmatter
\end{document}
